\documentclass[12pt]{scrartcl}
\usepackage[hagavi]{azzam}

% Modulo dan Keterbagian - Soal Latihan
% 72 soal: 17 Keterbagian, 11 FPB dan KPK, 29 Modulo, 15 Prima dan Pangkat.
% Sumber tiap soal dimunculkan di PDF sebagai \item (Sumber); komentar menyimpan detail lengkap.
% Versi gabungan Materi+Latihan tersedia sebagai Modulo-Keterbagian_Modul.tex.
% Soal pilihan ganda asli disajikan tanpa opsi jawaban sebagai soal isian,
% kecuali soal konsep nomor 1 yang memang membutuhkan opsinya.

\title{Modulo dan Keterbagian - Soal Latihan}
\date{12 Oktober 2026}

\begin{document}
\maketitle

\renewcommand*\contentsname{Daftar Isi}
\tableofcontents

\newpage

\section{Keterbagian}
\begin{enumerate}

	% Source: OSK SMA 2004 No. 4 (Pilihan Ganda)
	% Jawaban: B
	\item (OSK 2004)
	Untuk $a$ dan $b$ bilangan bulat dengan $a \ne 0$, notasi $a \mid b$ menyatakan
	``$a$ membagi $b$''. Pernyataan berikut yang \textbf{salah} adalah \dots
	\begin{enumerate}
		\item Jika $a \mid b$ dan $a \mid c$, maka $a \mid (bc)$.
		\item Jika $a \mid c$ dan $b \mid c$, maka $(ab) \mid c$.
		\item Jika $a \mid b$ dan $a \mid c$, maka $a \mid (b + c)$.
		\item Untuk setiap bilangan bulat $a \ne 0$ berlaku $a \mid 0$.
		\item Jika $a \mid b$, maka $a \mid (bc)$ untuk setiap bilangan bulat $c$.
	\end{enumerate}

	% Source: OSK SMA 2003 No. 1 (Pilihan Ganda)
	% Jawaban: 0
	\item (OSK 2003)
	Ada berapa banyak di antara bilangan-bilangan $20000002$, $20011002$,
	$20022002$, $20033002$ yang habis dibagi $9$?

	% Source: AIME 1984 (diadaptasi dari Modul Pra-OSK Keterbagian)
	% Jawaban: 592
	\item (AIME 1984)
	Bilangan asli $n$ adalah bilangan asli terkecil yang habis dibagi $15$
	dan semua digitnya sama dengan $8$ atau $0$. Tentukan nilai dari $\frac{n}{15}$.

	% Source: OSK SMA 2008 No. 5 (Pilihan Ganda)
	% Jawaban: 2
	\item (OSK 2008)
	Jumlah empat bilangan asli berturutan senantiasa habis dibagi $p$.
	Maka nilai $p$ terbesar adalah \dots

	% Source: OSP SMA 2015 No. 2 (Isian Singkat)
	% Jawaban: 4
	\item (OSP 2015)
	Banyaknya bilangan bulat $n$ sehingga $n + 1$ merupakan faktor dari $n^2 + 1$ adalah \dots

	% Source: OSK SMA 2013 No. 16 (Isian Singkat)
	% Jawaban: 64
	\item (OSK 2013)
	Diberikan himpunan $S = \left\{x \in \mathbb{Z} \;\middle|\; \frac{x^2 - 2x + 7}{2x - 1} \in \mathbb{Z}\right\}$.
	Banyaknya himpunan bagian dari $S$ adalah \dots

	% Source: OSK SMA 2014 No. 11 (Isian Singkat)
	% Jawaban: 595
	\item (OSK 2014)
	Bilangan-bilangan $1111$, $5276$, $8251$, dan $9441$ bersisa sama jika dibagi $N$.
	Nilai $N$ terbesar yang memiliki sifat tersebut adalah \dots

	% Source: OSK SMA 2012 No. 3 (Isian Singkat)
	% Jawaban: 960
	\item (OSK 2012)
	Bilangan terbesar $x$ kurang dari $1000$ sehingga terdapat tepat dua bilangan asli $n$
	dengan $\frac{n^2 + x}{n + 1}$ merupakan bilangan asli adalah \dots

	% Source: Finlandia MO 2001 (diadaptasi dari Modul Pra-OSK Keterbagian)
	% Jawaban: n = 6, 9, 2001
	\item (Finlandia MO 2001)
	Tentukan semua bilangan asli $n$ sedemikian hingga $n^2 + 2$ habis membagi $2001n + 2$.

	% Source: OSP SMA 2023 No. 2 (Isian Singkat)
	% Jawaban: 74
	\item (OSP 2023)
	Tentukan banyaknya bilangan asli $n \le 221$ sehingga
	\[ \frac{1^2 + 2^2 + \dots + n^2}{1 + 2 + \dots + n} \]
	merupakan bilangan bulat.

	% Source: Afrika Selatan MO 2021 (Modul Pra-OSK Keterbagian)
	% Jawaban: terkecil 17171, terbesar 98989
	\item (Afrika Selatan MO 2021)
	Tentukan bilangan asli terkecil dan terbesar dengan representasi digit
	yang berbentuk $ababa$ ($a \ne 0$) dan habis dibagi $11$.

	% Source: Kualifikasi Chile MO 2017 (Modul Pra-OSK Keterbagian)
	% Jawaban: d = 5 atau d = 8
	\item (Chile MO 2017)
	Dengan menggunakan dua digit berbeda $2$ dan $d$,
	bilangan $90$ digit $m = 2d22d222d2222d\ldots$ (blok berikutnya terdiri dari
	lima digit $2$ diikuti $d$, dan seterusnya) habis dibagi oleh $9$.
	Tentukan semua nilai yang mungkin dari $d$.

	% Source: Modul Pra-OSK Keterbagian
	% Jawaban: n = 2, 5, 8, 10, 15, 20, 26, 35, 40, 50, 65, 80, 90, 140, 170, 215, 290, 440, 890
	\item (Pra-OSK Keterbagian)
	Carilah semua bilangan asli $n$ sehingga $n + 10 \mid n^3 + 100$.

	% Source: India PRMO 2014 (Modul Pra-OSK Keterbagian)
	% Jawaban: 1
	\item (India PRMO 2014)
	Banyak bilangan asli $n$ di antara $1$ sampai $2014$ sehingga
	$\frac{8n}{9999 - n}$ merupakan bilangan bulat adalah \dots

	% Source: Modul Pra-OSK Keterbagian
	% Jawaban: n = k^2, k^2 + k, atau k^2 + 2k untuk suatu bilangan asli k
	\item (Pra-OSK Keterbagian)
	Tentukan semua bilangan asli $n$ sedemikian hingga $\lfloor \sqrt{n} \rfloor$ membagi $n$.

	% Source: Modul Pra-OSK Keterbagian
	% Jawaban: 1152
	\item (Pra-OSK Keterbagian)
	Diberikan sebuah barisan sepuluh bilangan asli naik tegas sedemikian hingga
	setiap bilangan membagi bilangan berikutnya.
	Jika jumlah kesepuluh bilangan tersebut adalah $1918$,
	maka semua kemungkinan bilangan kesepuluh dari barisan tersebut adalah \dots

	% Source: El Salvador MO 2020 (Modul Pra-OSK Keterbagian, notasi dirapikan)
	% Jawaban: 505
	\item (El Salvador MO 2020)
	Perhatikan bilangan $x_n = \overline{ab\,aab\,aaab\,\cdots\,\underbrace{a\ldots a}_{n}b}$,
	yaitu blok ke-$k$ terdiri dari $k$ digit $a$ diikuti satu digit $b$,
	dengan $a \ne 0$ dan $b$ keduanya digit.
	Tentukan banyak bilangan asli $n$ dengan $1 \le n \le 2020$ sedemikian hingga
	$x_n$ habis dibagi oleh $11$ untuk semua nilai $a$ dan $b$.

\end{enumerate}

\section{FPB dan KPK}
\begin{enumerate}[resume]

	% Source: Afrika Selatan MO 2006 (diadaptasi dari Modul Pra-OSK Keterbagian)
	% Jawaban: 70/37
	\item (Afrika Selatan MO 2006)
	Sederhanakan pecahan $\dfrac{2121212121210}{1121212121211}$ ke bentuk paling sederhana.

	% Source: OSK SMA 2024 No. 4 (Isian Singkat)
	% Jawaban: 214
	\item (OSK 2024)
	Misalkan $a$ dan $b$ bilangan bulat positif yang tidak memiliki faktor persekutuan selain $1$.
	Jika berlaku
	\[ \frac{1 + 2 + 3 + \dots + 104}{3 + 4 + 5 + \dots + 106} = \frac{a}{b}, \]
	maka nilai dari $a + b$ adalah \dots

	% Source: OSP SMA 2026 No. 5 (Isian Singkat)
	% Jawaban: 9
	\item (OSP 2026)
	Banyaknya pasangan terurut bilangan asli $(a, b)$ sehingga $a$ membagi $b^2$,
	$b$ membagi $a^2$, dan $ab = 10^{10}$ adalah \dots

	% Source: Modul Pra-OSK Keterbagian
	% Jawaban: 1 dan 12167
	\item (Pra-OSK Keterbagian)
	Hitunglah $FPB(133, 198)$ dan $FPB(6436343, 12167)$.

	% Source: Modul Pra-OSK Keterbagian
	% Jawaban: 32768
	\item (Pra-OSK Keterbagian)
	Berapa banyak pasangan bilangan asli $(x, y)$ yang memenuhi
	$FPB(x, y) = 5!$ dan $KPK(x, y) = 50!$?

	% Source: Modul Pra-OSK Keterbagian
	% Jawaban: 179
	\item (Pra-OSK Keterbagian)
	Definisikan $f(n) = FPB(n! + 1, (n + 1)!)$.
	Tentukan nilai dari $f(1) + f(2) + f(3) + \dots + f(30)$.

	% Source: OSK SMA 2021 Kemampuan Lanjut (Modul Pra-OSK Keterbagian)
	% Jawaban: 91
	\item (OSK 2021 Lanjut)
	Bilangan asli ganjil $b$ terbesar sedemikian hingga barisan bilangan asli
	$a_n = n^2 + 19n + b$ memenuhi
	\[ FPB(a_n, a_{n + 1}) = FPB(a_{n + 2}, a_{n + 1}) \]
	untuk setiap bilangan asli $n$ adalah \dots

	% Source: Modul Pra-OSK Keterbagian
	\item (Pra-OSK Keterbagian)
	Misalkan $m$ dan $n$ adalah bilangan asli yang memenuhi
	\[ FPB(m, n) + KPK(m, n) = m + n. \]
	Buktikan bahwa salah satu dari kedua bilangan itu membagi bilangan yang lainnya.

	% Source: Modul Pra-OSK Keterbagian
	\item (Pra-OSK Keterbagian)
	Diberikan bilangan asli $a$ dan $b$ yang memenuhi
	\[ \frac{a + 1}{b} + \frac{b + 1}{a} \]
	adalah bilangan bulat. Tunjukkan bahwa $FPB(a, b) \le \sqrt{a + b}$.

	% Source: Modul Pra-OSK Keterbagian
	\item (Pra-OSK Keterbagian)
	Misalkan $x$, $y$, $z$ adalah bilangan asli yang memenuhi
	\[ \frac{1}{x} - \frac{1}{y} = \frac{1}{z}. \]
	Misalkan $h = FPB(x, y, z)$.
	Tunjukkan bahwa $hxyz$ dan $h(y - x)$ keduanya merupakan bilangan kuadrat.

	% Source: Modul Pra-OSK Keterbagian
	% Jawaban: (t, t, t) dan semua permutasi dari (t, 2t, 3t) untuk t bilangan asli
	\item (Pra-OSK Keterbagian)
	Carilah semua tripel bilangan asli $(a, b, c)$ sedemikian hingga
	$a^3 + b^3 + c^3$ habis dibagi $a^2b$, $b^2c$, dan $c^2a$.

\end{enumerate}

\section{Modulo}
\begin{enumerate}[resume]

	% Source: OSK SMA 2009 No. 18 (Isian Singkat)
	% Jawaban: 6
	\item (OSK 2009)
	Jika $10^{999999999}$ dibagi oleh $7$, maka sisanya adalah \dots

	% Source: OSK SMA 2005 No. 8 (Pilihan Ganda, disajikan sebagai isian)
	% Jawaban: 6
	\item (OSK 2005)
	Tentukan digit terakhir dari $8^{8^{8^8}}$.

	% Source: OSK SMA 2003 No. 9 (Pilihan Ganda)
	% Jawaban: 13
	\item (OSK 2003)
	Misalkan $N$ adalah bilangan bulat terkecil yang bersifat bersisa $2$ jika dibagi $5$,
	bersisa $3$ jika dibagi $7$, dan bersisa $4$ jika dibagi $9$.
	Berapakah hasil penjumlahan digit-digit dari $N$?

	% Source: OSK SMA 2021 No. 3 (Isian Singkat)
	% Jawaban: 21
	\item (OSK 2021)
	Diberikan bilangan bulat positif $A$ dan $B$ yang bila dibagi $5$ berturut-turut bersisa $2$ dan $3$.
	Sisa pembagian $A(A + 1) + 5B$ jika dibagi $25$ adalah \dots

	% Source: OSK SMA 2023 No. 8 (Isian Singkat)
	% Jawaban: 30
	\item (OSK 2023)
	Sisa pembagian dari $5^{2021} + 11^{2022}$ oleh $64$ adalah \dots

	% Source: OSP SMA 2002 No. 16 (Isian Singkat)
	% Jawaban: 7
	\item (OSP 2002)
	Berapakah sisa pembagian $43^{43^{43}}$ oleh $100$?

	% Source: OSP SMA 2011 No. 17 (Isian Singkat)
	% Jawaban: semua bilangan asli
	\item (OSP 2011)
	Bilangan asli $n$ yang memenuhi $(-2004)^n - 1900^n + 25^n - 121^n$ habis dibagi $2000$ adalah \dots

	% Source: Canadian MO Qualification Repechage 2016 (diadaptasi dari Modul Pra-OSK Modulo)
	% Jawaban: tidak ada
	\item (Canadian MO 2016)
	Tentukan semua bilangan asli $n$ sehingga $11 \mid (3^n + 4^n)$.

	% Source: OSK SMA 2021 No. 3 (Isian Singkat)
	% Jawaban: 17
	\item (OSK 2021)
	Diketahui dua digit terakhir dari $a^{777}$ adalah $77$.
	Maka dua digit terakhir dari $a$ adalah \dots

	% Source: Modul Pra-OSK Modulo
	% Jawaban: 43
	\item (Pra-OSK Modulo)
	Tentukan dua digit terakhir dari $7^{707}$.

	% Source: Modul Pra-OSK Modulo
	% Jawaban: 9 dan 1
	\item (Pra-OSK Modulo)
	Carilah digit terakhir dari $13^{14}$ dan $13^{14^{15}}$.

	% Source: Modul Pra-OSK Modulo
	% Jawaban: 625
	\item (Pra-OSK Modulo)
	Tentukan sisa pembagian saat $1 \times 3 \times 5 \times \dots \times 2023$ dibagi oleh $1000$.

	% Source: Modul Pra-OSK Modulo
	% Jawaban: 76 dan 77
	\item (Pra-OSK Modulo)
	Carilah dua digit terakhir dari $1032^{1032}$ dan $17^{17^{17}}$.

	% Source: Argentina MO 2021 Level 1 (Modul Pra-OSK Modulo)
	% Jawaban: 9
	\item (Argentina MO 2021)
	Mikha menulis barisan aritmetika yang dimulai dari $1$ dengan selisih $3$, yakni
	\[ 1, 4, 7, 10, 13, 16, \ldots \]
	Kemudian, Fakhru menandai semua bilangan yang ditulis Mikha yang lebih dari $10$
	dan kurang dari $100000$ dan seluruh digitnya sama.
	Tentukan banyak bilangan yang ditandai oleh Fakhru.

	% Source: Modul Pra-OSK Modulo
	\item (Pra-OSK Modulo)
	Misalkan $m$, $a$, $b$ adalah bilangan bulat dengan $a$ dan $b$ relatif prima dengan $m$.
	Bilangan asli $x$ dan $y$ memenuhi $a^x \equiv b^x \pmod{m}$ dan $a^y \equiv b^y \pmod{m}$.
	Tunjukkan bahwa $a^{FPB(x, y)} \equiv b^{FPB(x, y)} \pmod{m}$.

	% Source: Modul Pra-OSK Modulo (pernyataan diperbaiki: p ganjil)
	\item (Pra-OSK Modulo)
	Misalkan $p$ adalah bilangan prima ganjil.
	Tunjukkan bahwa terdapat bilangan asli $x$ sedemikian hingga $x^2 + 1$
	habis dibagi oleh $p$ jika dan hanya jika $p \equiv 1 \pmod{4}$.

	% Source: Modul Pra-OSK Modulo (pernyataan diperbaiki: p - 1 digit)
	\item (Pra-OSK Modulo)
	Misalkan $p \ge 7$ adalah bilangan prima.
	Tunjukkan bahwa bilangan $11\ldots11$ (dengan $p - 1$ digit $1$) habis dibagi oleh $p$.

	% Source: Modul Pra-OSK Modulo
	\item (Pra-OSK Modulo)
	Misalkan $p$ adalah bilangan prima yang berbentuk $3k + 2$ yang membagi $a^2 + ab + b^2$
	untuk suatu bilangan bulat $a$ dan $b$.
	Buktikan bahwa $a$ dan $b$ keduanya habis dibagi oleh $p$.

	% Source: Modul Pra-OSK Modulo
	\item (Pra-OSK Modulo)
	Misalkan $n > 1$ adalah bilangan ganjil. Tunjukkan bahwa $n$ tidak membagi $3^n + 1$.

	% Source: Modul Pra-OSK Modulo
	\item (Pra-OSK Modulo)
	Misalkan $a$ adalah bilangan ganjil.
	Buktikan bahwa $a^{2^n} + 2^{2^n}$ dan $a^{2^m} + 2^{2^m}$ relatif prima saat $m \ne n$.

	% Source: Modul Pra-OSK Modulo (pernyataan diperbaiki: n > 1)
	\item (Pra-OSK Modulo)
	Misalkan $n > 1$ adalah bilangan asli dan $p$ bilangan prima sedemikian hingga
	$n$ membagi $p - 1$ dan $p$ membagi $n^3 - 1$.
	Tunjukkan bahwa $4p - 3$ adalah bilangan kuadrat.

	% Source: Modul Pra-OSK Modulo
	\item (Pra-OSK Modulo)
	Jika $p$ adalah bilangan prima ganjil dan $a$, $b$ adalah dua bilangan asli yang relatif prima,
	maka tunjukkan bahwa $FPB\left(a + b, \dfrac{a^p + b^p}{a + b}\right) = 1$ atau $p$.

	% Source: Modul Pra-OSK Modulo
	% Jawaban: (p - 1)p^2/2
	\item (Pra-OSK Modulo)
	Misalkan $p$ adalah bilangan prima ganjil.
	Untuk setiap $i = 1, 2, \ldots, p - 1$, notasikan $r_i$ sebagai sisa pembagian
	ketika $i^p$ dibagi dengan $p^2$.
	Carilah nilai dari $r_1 + r_2 + \dots + r_{p - 1}$.

	% Source: Modul Pra-OSK Modulo
	\item (Pra-OSK Modulo)
	Buktikan bahwa untuk setiap bilangan genap positif $n$,
	$n^2 - 1$ membagi $2^{n!} - 1$.

	% Source: Modul Pra-OSK Modulo
	\item (Pra-OSK Modulo)
	Misalkan $a$ dan $b > 2$ adalah bilangan asli.
	Tunjukkan bahwa $2^a + 1$ tidak habis dibagi oleh $2^b - 1$.

	% Source: Modul Pra-OSK Modulo
	\item (Pra-OSK Modulo)
	Misalkan $n > 2$ adalah bilangan asli.
	Tunjukkan bahwa $n^{n^{n^n}} - n^{n^n}$ habis dibagi oleh $1989$.

	% Source: Modul Pra-OSK Modulo
	\item (Pra-OSK Modulo)
	Misalkan $a$ dan $b$ adalah dua bilangan asli sedemikian hingga $b^n + n$
	merupakan kelipatan dari $a^n + n$ untuk setiap bilangan asli $n$.
	Tunjukkan bahwa $a = b$.

	% Source: Modul Pra-OSK Modulo
	\item (Pra-OSK Modulo)
	Misalkan $x$, $a$, $b$ adalah bilangan asli.
	Tunjukkan bahwa $FPB(x^a - 1, x^b - 1) = x^{FPB(a, b)} - 1$.

	% Source: Modul Pra-OSK Modulo
	\item (Pra-OSK Modulo)
	Misalkan $a$ adalah bilangan asli dan $p$ bilangan prima.
	Tunjukkan bahwa terdapat suatu bilangan asli $m$ yang memenuhi
	\[ m^{m^m} \equiv a \pmod{p}. \]

\end{enumerate}

\section{Bilangan Prima dan Pangkat}
\begin{enumerate}[resume]

	% Source: OSK SMA 2011 (diadaptasi dari Modul Pra-OSK Modulo)
	% Jawaban: p = 3
	\item (OSK 2011)
	Tentukan semua bilangan asli $p$ sedemikian hingga $p$, $p + 8$, dan $p + 16$
	semuanya merupakan bilangan prima.

	% Source: Modul Pra-OSK Modulo
	% Jawaban: (a, b) = (1, 1)
	\item (Pra-OSK Modulo)
	Tentukan semua bilangan asli $a$ dan $b$ sedemikian hingga $a^4 + 4b^4$
	merupakan bilangan prima.

	% Source: Modul Pra-OSK Modulo
	% Jawaban: (m, n) = (1, 3)
	\item (Pra-OSK Modulo)
	Carilah semua bilangan asli $m$ dan $n$ sedemikian hingga $6^m + 2^n + 2$
	adalah bilangan kuadrat.

	% Source: Modul Pra-OSK Modulo
	% Jawaban: n = 2
	\item (Pra-OSK Modulo)
	Tentukan semua bilangan asli $n$ sedemikian hingga $3n - 4$, $4n - 5$,
	dan $5n - 3$ semuanya merupakan bilangan prima.

	% Source: Modul Pra-OSK Modulo
	\item (Pra-OSK Modulo)
	Buktikan bahwa setiap bilangan prima yang lebih besar dari $3$
	berbentuk $6k \pm 1$ untuk suatu bilangan asli $k$.

	% Source: Modul Pra-OSK Modulo (soal yang sama sudah dibahas sebagai contoh di Materi)
	\item (Pra-OSK Modulo)
	Buktikan bahwa untuk setiap bilangan prima $p > 3$ berlaku $24 \mid p^2 - 1$.

	% Source: Modul Pra-OSK Modulo
	% Jawaban: 0
	\item (Pra-OSK Modulo)
	Banyak bilangan di barisan $11, 111, 1111, 11111, \ldots, 11\ldots11$
	dengan bilangan terakhir terdiri dari $2022$ digit $1$,
	yang merupakan bilangan kuadrat adalah \dots

	% Source: Modul Pra-OSK Modulo
	% Jawaban: n = 5
	\item (Pra-OSK Modulo)
	Tentukan semua bilangan asli $n$ sehingga $n! + 5$ merupakan bilangan kubik.

	% Source: Argentina MO 2021 Level 1 (Modul Pra-OSK Modulo)
	% Jawaban: (p, q) = (7, 3), (13, 3), (37, 3), (97, 3)
	\item (Argentina MO 2021)
	Tentukan semua bilangan prima $p$ dan $q$ yang lebih dari $1$ dan kurang dari $100$
	sedemikian hingga lima bilangan berikut:
	\[ p + 6,\quad p + 10,\quad q + 4,\quad q + 10,\quad p + q + 1 \]
	semuanya merupakan bilangan prima.

	% Source: Modul Pra-OSK Modulo
	\item (Pra-OSK Modulo)
	Jika $p$ adalah bilangan prima ganjil dan
	$\dfrac{a}{b} = \dfrac{1}{1} + \dfrac{1}{2} + \dots + \dfrac{1}{p - 1}$
	dalam bentuk paling sederhana,
	buktikan bahwa $p \mid a$.

	% Source: Modul Pra-OSK Modulo (pernyataan diperbaiki: memuat kata prima)
	\item (Pra-OSK Modulo)
	Buktikan bahwa jika $a^n + 1$ merupakan bilangan prima dengan $a > 1$,
	maka $a$ genap dan $n = 2^k$ untuk suatu bilangan cacah $k$.
	(Bilangan prima berbentuk $2^{2^k} + 1$ disebut bilangan prima Fermat.)

	% Source: Modul Pra-OSK Modulo (pernyataan diperbaiki: memuat kata prima)
	\item (Pra-OSK Modulo)
	Buktikan bahwa jika $a^n - 1$ merupakan bilangan prima dengan $a > 1$,
	maka $a = 2$ dan $n$ prima.
	(Bilangan prima berbentuk $2^n - 1$ disebut bilangan prima Mersenne.)

	% Source: Modul Pra-OSK Modulo
	\item (Pra-OSK Modulo)
	Buktikan bahwa untuk setiap bilangan asli $n$,
	terdapat $n$ buah bilangan asli berurutan yang tidak ada satupun
	di antaranya yang merupakan bilangan berpangkat sempurna.

	% Source: Modul Pra-OSK Modulo (redaksi diperbaiki)
	% Jawaban: ya, selalu ada
	\item (Pra-OSK Modulo)
	Apakah untuk setiap himpunan bilangan asli $\{a_1, a_2, \ldots, a_k\}$
	terdapat suatu bilangan asli $b$ sedemikian hingga setiap anggota himpunan
	$\{ba_1, ba_2, \ldots, ba_k\}$ merupakan bilangan berpangkat sempurna?

	% Source: OSP SMA 2010 (Modul Pra-OSK Modulo)
	% Jawaban: p = 13
	\item (OSP 2010)
	Tentukan semua bilangan prima $p$ sehingga $p^2 + 7^3$ merupakan bilangan kubik.

\end{enumerate}

\end{document}
